% !TeX root = ../main.tex
% =========================================================================
% 4. OBJETIVOS (general: qué + para qué + cómo | específicos: Bloom)
% =========================================================================
\section{Objetivos}

\subsection{Objetivo General}
Desarrollar un agente inteligente para la atención comercial y cotización de materiales eléctricos en Guayaquil, mediante la integración de modelos generativos y recuperación relacional determinista sobre catálogos técnicos para garantizar cotizaciones fácticas en la distribuidora Power Electrimk.

\subsection{Objetivos Específicos}

De conformidad con los niveles taxonómicos de Bloom y las directrices metodológicas de titulación en ingeniería computacional de la Universidad Politécnica Salesiana, se establecen los siguientes objetivos específicos:

\begin{itemize}
    \item \textbf{Analizar} los requerimientos operativos de atención comercial y la estructura de los catálogos técnicos multimarca en la distribuidora Power Electrimk, determinando los atributos esenciales de inventario y las especificaciones de ingeniería requeridas para el agente de ventas.
    \item \textbf{Implementar} un prototipo de software en arquitectura ligera sobre .NET 8 que integre de forma desacoplada la API oficial de WhatsApp Cloud, el modelo fundacional Gemini y un motor de base de datos embebido SQLite con registro en modo Write-Ahead Logging para consultas deterministas y persistencia de memoria de sesión comercial.
    \item \textbf{Evaluar} el desempeño computacional y la efectividad operativa del prototipo mediante pruebas experimentales controladas que contrasten el agente con recuperación local frente al modelo sin recuperación, cuantificando la exactitud de precios y referencias técnicas, la tasa de mitigación de alucinaciones, la latencia en percentil P95 y el consumo de recursos en la nube.
\end{itemize}
